\documentclass[10pt,a4paper]{amsart}
\usepackage{amsmath,amssymb,amsthm,graphicx}
% Retained preamble for the Stacks Project chapter divisors.tex (Divisors).
% The chapter file pinned-original-source.tex carries only \input{preamble} on its first line, so no
% preamble text lives inside the chapter itself.  The Stacks Project uses one single shared preamble for
% every chapter; the snapshot stages it verbatim in sources/stacks-chow/preamble.tex.  The block below is
% copied mechanically, line for line, from that file's lines 14-18 (the Xy-pic setup) and lines 154-217
% (the native theorem declarations together with the mathematical macros).  Nothing was added, removed or
% redefined.  The wrapper class and the amsmath/amssymb/amsthm/graphicx base packages are supplied by the
% assembler, so the class/\externaldocument part of the shared preamble is not repeated here.
% For commutative diagrams we use Xy-pic
\usepackage[all]{xy}

% We use 2cell for 2-commutative diagrams.
\xyoption{2cell}
% Theorem environments.
%
\theoremstyle{plain}
\newtheorem{theorem}[subsection]{Theorem}
\newtheorem{proposition}[subsection]{Proposition}
\newtheorem{lemma}[subsection]{Lemma}

\theoremstyle{definition}
\newtheorem{definition}[subsection]{Definition}
\newtheorem{example}[subsection]{Example}
\newtheorem{exercise}[subsection]{Exercise}
\newtheorem{situation}[subsection]{Situation}

\theoremstyle{remark}
\newtheorem{remark}[subsection]{Remark}
\newtheorem{remarks}[subsection]{Remarks}

\numberwithin{equation}{subsection}

% Macros
%
\def\lim{\mathop{\mathrm{lim}}\nolimits}
\def\colim{\mathop{\mathrm{colim}}\nolimits}
\def\Spec{\mathop{\mathrm{Spec}}}
\def\Hom{\mathop{\mathrm{Hom}}\nolimits}
\def\Ext{\mathop{\mathrm{Ext}}\nolimits}
\def\SheafHom{\mathop{\mathcal{H}\!\mathit{om}}\nolimits}
\def\SheafExt{\mathop{\mathcal{E}\!\mathit{xt}}\nolimits}
\def\Sch{\mathit{Sch}}
\def\Mor{\mathop{\mathrm{Mor}}\nolimits}
\def\Ob{\mathop{\mathrm{Ob}}\nolimits}
\def\Sh{\mathop{\mathit{Sh}}\nolimits}
\def\NL{\mathop{N\!L}\nolimits}
\def\CH{\mathop{\mathrm{CH}}\nolimits}
\def\proetale{{pro\text{-}\acute{e}tale}}
\def\etale{{\acute{e}tale}}
\def\QCoh{\mathit{QCoh}}
\def\Ker{\mathop{\mathrm{Ker}}}
\def\Im{\mathop{\mathrm{Im}}}
\def\Coker{\mathop{\mathrm{Coker}}}
\def\Coim{\mathop{\mathrm{Coim}}}

% Boxtimes
%
\DeclareMathSymbol{\boxtimes}{\mathbin}{AMSa}{"02}

%
% Macros for moduli stacks/spaces
%
\def\QCohstack{\mathcal{QC}\!\mathit{oh}}
\def\Cohstack{\mathcal{C}\!\mathit{oh}}
\def\Spacesstack{\mathcal{S}\!\mathit{paces}}
\def\Quotfunctor{\mathrm{Quot}}
\def\Hilbfunctor{\mathrm{Hilb}}
\def\Curvesstack{\mathcal{C}\!\mathit{urves}}
\def\Polarizedstack{\mathcal{P}\!\mathit{olarized}}
\def\Complexesstack{\mathcal{C}\!\mathit{omplexes}}
% \Pic is the operator that assigns to X its picard group, usage \Pic(X)
% \Picardstack_{X/B} denotes the Picard stack of X over B
% \Picardfunctor_{X/B} denotes the Picard functor of X over B
\def\Pic{\mathop{\mathrm{Pic}}\nolimits}
\def\Picardstack{\mathcal{P}\!\mathit{ic}}
\def\Picardfunctor{\mathrm{Pic}}
\def\Deformationcategory{\mathcal{D}\!\mathit{ef}}

% --- appended 2026-10-12 (Chief): the project own definitions of the reference/slogan/history
% display environments, copied mechanically from the shared Stacks preamble lines 8-12
% (\usepackage{verbatim} plus the three \newenvironment lines, which the chapter preamble keeps
% verbatim from the project preamble). Nothing added, removed or redefined: these environments
% carry no mathematical content and the project itself typesets them as comments. Without these
% three lines a counted statement whose own text uses them fails with Environment slogan undefined.
\usepackage{verbatim}
\newenvironment{reference}{\comment}{\endcomment}
\newenvironment{slogan}{\comment}{\endcomment}
\newenvironment{history}{\comment}{\endcomment}
\usepackage{cleveref}
\crefname{originallemma}{original lemma}{original lemmas}
\Crefname{originallemma}{Original lemma}{Original lemmas}
\makeatletter
\newlabel{algebra-lemma-criterion-no-embedded-primes}{{algebra-lemma-criterion-no-embedded-primes}{source}{}{sourcepointer-1}{}}
\newlabel{algebra-lemma-criterion-no-embedded-primes@cref}{{[originallemma][0][]algebra-lemma-criterion-no-embedded-primes}{[1][0][]source}{}{}{}}
\makeatother
\crefname{originallemma}{original lemma}{original lemmas}
\Crefname{originallemma}{Original lemma}{Original lemmas}
\makeatletter
\newlabel{algebra-lemma-tensor-fields-CM}{{algebra-lemma-tensor-fields-CM}{source}{}{sourcepointer-2}{}}
\newlabel{algebra-lemma-tensor-fields-CM@cref}{{[originallemma][0][]algebra-lemma-tensor-fields-CM}{[1][0][]source}{}{}{}}
\makeatother
\crefname{originalproposition}{original proposition}{original propositions}
\Crefname{originalproposition}{Original proposition}{Original propositions}
\makeatletter
\newlabel{algebra-proposition-minimal-primes-associated-primes}{{algebra-proposition-minimal-primes-associated-primes}{source}{}{sourcepointer-3}{}}
\newlabel{algebra-proposition-minimal-primes-associated-primes@cref}{{[originalproposition][0][]algebra-proposition-minimal-primes-associated-primes}{[1][0][]source}{}{}{}}
\makeatother
\crefname{originallemma}{original lemma}{original lemmas}
\Crefname{originallemma}{Original lemma}{Original lemmas}
\makeatletter
\newlabel{lemma-bourbaki}{{05DB}{source}{}{sourcepointer-4}{}}
\newlabel{lemma-bourbaki@cref}{{[originallemma][0][]05DB}{[1][0][]source}{}{}{}}
\makeatother
\crefname{originallemma}{original lemma}{original lemmas}
\Crefname{originallemma}{Original lemma}{Original lemmas}
\makeatletter
\newlabel{morphisms-lemma-finite-type-noetherian}{{morphisms-lemma-finite-type-noetherian}{source}{}{sourcepointer-5}{}}
\newlabel{morphisms-lemma-finite-type-noetherian@cref}{{[originallemma][0][]morphisms-lemma-finite-type-noetherian}{[1][0][]source}{}{}{}}
\makeatother
\title{Relative associated points under base change are the associated points upstairs together with generic points of the residue field tensor products}
\author{The Stacks Project Authors}
\date{stacks-project:divisors.tex:local-snapshot-20261012; GNU FDL 1.2 or later}
\begin{document}
\maketitle
\noindent\textit{Editorial note.} All delivered mathematical text is the author's own, copied from the pinned Stacks Project chapter file pinned-original-source.tex; no statement or proof step was rewritten, reordered or removed.\par\medskip
\noindent\textit{Editorial note.} Following the multi-claim convention for this source, source\_url is the source-level URL https://stacks.math.columbia.edu/ and the source label lemma-base-change-relative-assassin is carried in claim\_key and here rather than in a per-claim tag URL.\par\medskip
\noindent\textit{Editorial note.} The preamble is the chapter's shared Stacks Project preamble, retained as sources/stacks-divisors/retained-preamble.tex: a line-for-line excerpt of the shared Stacks preamble staged with this snapshot. No macro was invented or redefined, and this card adds no preamble text of its own.\par\medskip
\noindent\textit{Editorial note.} Statement span 1119-1138 and proof span 1140-1167 come from the halfA pool card operations/pool-cards/stacks-divisors-halfA.json and were re-read in the pinned chapter file before this card was written.\par\medskip
\noindent\textit{Editorial note.} The statement contains an xymatrix base change diagram; it is copied verbatim and the retained chapter preamble loads Xy-pic.\par\medskip
\noindent\textit{Editorial note.} The notation Ass\_\{X/S\}(F) for relative associated points is the author's and is preserved.\par\medskip
\noindent\textit{Editorial note.} 5 cross-references occur inside the counted spans; they are listed in external\_references as navigation pointers only.\par\medskip
\noindent\textit{Allowed standard prerequisites (editorial).}\begin{itemize}
\item The same-chapter decomposition of the associated points of a quasi-coherent module on a scheme over a locally Noetherian base into the associated points of the fibres and the generic points of the fibres. Cited at proof lines 1146-1149. It is the tool that lets the associated points upstairs be written as a union indexed by the associated points of the fibre downstairs.
\item Fibres of a locally finite type morphism over the residue field of a point are locally Noetherian, cited from the morphisms chapter. Cited at proof lines 1144-1146. Makes the cited decomposition applicable to the fibre at the chosen point.
\item The residue field tensor product of a field extension is Noetherian and Cohen-Macaulay, and in such a ring the associated primes are exactly the minimal primes, cited from the algebra chapter. Cited at proof lines 1160-1167. Shows that the associated points of the fibre of the base changed morphism over the chosen point are precisely its generic points, which is the remaining step of the argument.
\item Flatness of the morphism between the two fibres, which is a base change of residue field extensions. Used at proof lines 1140-1143. Identifies the pullback of the module on the fibre with the module on the base changed fibre.
\end{itemize}\medskip
\ifdefined\hypertarget\hypertarget{sourcepointer-1}{}\fi
\noindent\textit{Source reference.} Original lemma algebra-lemma-criterion-no-embedded-primes: sources/stacks-divisors/pinned-original-source.tex, source lines 422--422. Same-chapter reference used in the retained author text at source line(s) 422, 1165. It is supplied only as a navigation pointer so that the retained text compiles; the counted item delivers the author's own statement and the author's complete proof, and no part of the referenced item is copied in or substituted for the counted argument.\par\medskip
\ifdefined\hypertarget\hypertarget{sourcepointer-2}{}\fi
\noindent\textit{Source reference.} Original lemma algebra-lemma-tensor-fields-CM: sources/stacks-divisors/pinned-original-source.tex, source lines 1160--1160. Same-chapter reference used in the retained author text at source line(s) 1160. It is supplied only as a navigation pointer so that the retained text compiles; the counted item delivers the author's own statement and the author's complete proof, and no part of the referenced item is copied in or substituted for the counted argument.\par\medskip
\ifdefined\hypertarget\hypertarget{sourcepointer-3}{}\fi
\noindent\textit{Source reference.} Original proposition algebra-proposition-minimal-primes-associated-primes: sources/stacks-divisors/pinned-original-source.tex, source lines 475--475. Same-chapter reference used in the retained author text at source line(s) 475, 1163. It is supplied only as a navigation pointer so that the retained text compiles; the counted item delivers the author's own statement and the author's complete proof, and no part of the referenced item is copied in or substituted for the counted argument.\par\medskip
\ifdefined\hypertarget\hypertarget{sourcepointer-4}{}\fi
\noindent\textit{Source reference.} Original lemma 05DB: focused-seeds/AGP/sources/stacks-snapshot-20261011/tags.txt, source lines 6949--6949. Same-chapter reference used in the retained author text at source line(s) 1149. It is supplied only as a navigation pointer so that the retained text compiles; the counted item delivers the author's own statement and the author's complete proof, and no part of the referenced item is copied in or substituted for the counted argument.\par\medskip
\ifdefined\hypertarget\hypertarget{sourcepointer-5}{}\fi
\noindent\textit{Source reference.} Original lemma morphisms-lemma-finite-type-noetherian: sources/stacks-divisors/pinned-original-source.tex, source lines 1147--1147. Same-chapter reference used in the retained author text at source line(s) 1147, 5702. It is supplied only as a navigation pointer so that the retained text compiles; the counted item delivers the author's own statement and the author's complete proof, and no part of the referenced item is copied in or substituted for the counted argument.\par\medskip
\begin{lemma}
\label{lemma-base-change-relative-assassin}
Let $f : X \to S$ be a morphism of schemes.
Let $\mathcal{F}$ be a quasi-coherent $\mathcal{O}_X$-module.
Let $g : S' \to S$ be a morphism of schemes.
Consider the base change diagram
$$
\xymatrix{
X' \ar[d] \ar[r]_{g'} & X \ar[d] \\
S' \ar[r]^g & S
}
$$
and set $\mathcal{F}' = (g')^*\mathcal{F}$. Let $x' \in X'$ be a point
with images $x \in X$, $s' \in S'$ and $s \in S$.
Assume $f$ locally of finite type.
Then $x' \in \text{Ass}_{X'/S'}(\mathcal{F}')$ if and only if
$x \in \text{Ass}_{X/S}(\mathcal{F})$ and $x'$ corresponds to
a generic point of an irreducible component of
$\Spec(\kappa(s') \otimes_{\kappa(s)} \kappa(x))$.
\end{lemma}

\begin{proof}
Consider the morphism $X'_{s'} \to X_s$ of fibres. As
$X_{s'} = X_s \times_{\Spec(\kappa(s))} \Spec(\kappa(s'))$
this is a flat morphism. Moreover $\mathcal{F}'_{s'}$ is the pullback
of $\mathcal{F}_s$ via this morphism. As $X_s$ is locally of finite
type over the Noetherian scheme $\Spec(\kappa(s))$ we have that
$X_s$ is locally Noetherian, see
Morphisms, Lemma \ref{morphisms-lemma-finite-type-noetherian}.
Thus we may apply
Lemma \ref{lemma-bourbaki}
and we see that
$$
\text{Ass}_{X'_{s'}}(\mathcal{F}'_{s'}) =
\bigcup\nolimits_{x \in \text{Ass}(\mathcal{F}_s)} \text{Ass}((X'_{s'})_x).
$$
Thus to prove the lemma it suffices to show that the associated points
of the fibre $(X'_{s'})_x$ of the morphism $X'_{s'} \to X_s$ over $x$
are its generic points. Note that
$(X'_{s'})_x = \Spec(\kappa(s') \otimes_{\kappa(s)} \kappa(x))$
as schemes. By
Algebra, Lemma \ref{algebra-lemma-tensor-fields-CM}
the ring $\kappa(s') \otimes_{\kappa(s)} \kappa(x)$ is a Noetherian
Cohen-Macaulay ring. Hence its associated primes are its minimal primes, see
Algebra, Proposition \ref{algebra-proposition-minimal-primes-associated-primes}
(minimal primes are associated) and
Algebra, Lemma \ref{algebra-lemma-criterion-no-embedded-primes}
(no embedded primes).
\end{proof}

\end{document}
